% Introduction générale
\chapter*{Introduction générale}
\addcontentsline{toc}{chapter}{Introduction générale}

L'intelligence artificielle et l'apprentissage automatique occupent une place croissante dans les systèmes de décision, industriels, médicaux, financiers ou réglementaires. La \textbf{combinaison de classifieurs}, ou apprentissage d'ensemble (\emph{ensemble learning}), s'impose comme un paradigme central : on construit un comité de modèles dont les prédictions sont agrégées pour améliorer la précision et la robustesse (bagging, boosting, stacking). Parallèlement, l'\textbf{explicabilité des modèles} est devenue une exigence majeure : les décisions automatisées sont scrutées par les utilisateurs, les régulateurs et la société ; comprendre \emph{pourquoi} un modèle prédit relève de la confiance, de la conformité et de l'éthique. Les méthodes d'explicabilité (intrinsèques ou post-hoc : LIME, SHAP, importance des variables) visent à rendre les systèmes moins opaques. Ce rapport se situe à l'intersection de ces deux thématiques : les systèmes d'ensemble complexifient le modèle global et renforcent l'effet «~boîte noire~», tout en offrant des gains de performance substantiels.

Les méthodes d'ensemble améliorent la précision au prix d'une complexité accrue (plusieurs modèles, fusion, dépendances), ce qui soulève une question centrale : \textit{comment concilier la puissance des systèmes combinant plusieurs classifieurs avec l'exigence d'explicabilité et d'interprétabilité des décisions ?} D'un côté, un classifieur unique peut être plus facile à expliquer ; de l'autre, un ensemble peut être plus performant mais plus difficile à interpréter. La problématique recouvre les dimensions techniques, architecturales et réglementaires (transparence, traçabilité).

Les objectifs sont de présenter les fondements de la combinaison de classifieurs (définitions, typologies, motivations théoriques, aspects architecturaux) ; de clarifier les concepts d'explicabilité et d'interprétabilité et les principales méthodes (intrinsèques et post-hoc : LIME, SHAP, PDP/ICE, explicabilité locale et globale) ; d'analyser les enjeux spécifiques aux ensembles (complexité, stratégies d'explication, compromis performance et compréhension) ; et de mettre en perspective les avantages et limites au regard de l'explicabilité.

Le document comporte trois parties : la \textbf{combinaison de classifieurs} (définition, typologies, aspects architecturaux, analyse comparative, avantages et limites) ; l'\textbf{explicabilité des modèles} (définition, interprétabilité, méthodes, enjeux dans les systèmes décisionnels, éthique et réglementation) ; la \textbf{combinaison et l'explicabilité} (boîte noire, stratégies, compromis performance et compréhension). La conclusion générale synthétise les apports et ouvre sur les perspectives.
